%%%%%%%%%%%%%%%%%%%%%%%%%%%%%%%%%%%%%%%%%%%%%%%%%
%%%%%%%%%%%%%%%%%% INFORMATION %%%%%%%%%%%%%%%%%%
%%%%%%%%%%%%%%%%%%%%%%%%%%%%%%%%%%%%%%%%%%%%%%%%%

% In those acronyms, "long" should explicit the use of letters if possible
% It will be displayed in the Acronyms section of the document
% Field "description" is what is displayed when hovering an acronym in the document

% If "long" is empty, the acronym will not appear in the list but will be available for tooltips
% This can be used for instance to have a plural version of an acronym to show in a tooltip without duplicating the entry

% In the document, use:
% * \acro{key}     -> Shows the short form, with the description as a tooltip     (e.g., "CNN")
% * \acrofull{key} -> Shows the description, followed by the short form           (e.g., "Convolutional Neural Network (CNN)")
% * \acroname{key} -> Shows the description only                                  (e.g., "Convolutional Neural Network")

% This file is shared by the three example documents (HDR, HDR demand, PhD)
% Entries are sorted alphabetically, which is not required but makes edition easier

%%%%%%%%%%%%%%%%%%%%%%%%%%%%%%%%%%%%%%%%%%%%%%%%%
%%%%%%%%%%%%%%%%%%%% ACRONYMS %%%%%%%%%%%%%%%%%%%
%%%%%%%%%%%%%%%%%%%%%%%%%%%%%%%%%%%%%%%%%%%%%%%%%

\newglossaryentry{ai}
{
    type        = {acronym},
    name        = {AI},
    long        = {\textbf{A}rtificial \textbf{I}ntelligence},
    description = {Artificial Intelligence}
}

% Example of a French counterpart of an existing acronym
% All entries of this file are listed in every document that prints a list of acronyms,
% so a language-specific entry is better declared with an empty "long": it then keeps
% its tooltip, without appearing in the list of the documents written in another language
\newglossaryentry{ai_fr}
{
    type        = {acronym},
    name        = {IA},
    long        = {},
    description = {Intelligence Artificielle}
}

\newglossaryentry{cnn}
{
    type        = {acronym},
    name        = {CNN},
    long        = {\textbf{C}onvolutional \textbf{N}eural \textbf{N}etwork},
    description = {Convolutional Neural Network}
}

% Example of a tooltip-only entry: "long" is empty, so it does not pollute the list of acronyms
% Here, this gives a plural form without duplicating the entry above
\newglossaryentry{cnn_plural}
{
    type        = {acronym},
    name        = {CNNs},
    long        = {},
    description = {Convolutional Neural Networks}
}

% Same as above: only used in the French example, hence not listed
\newglossaryentry{ed}
{
    type        = {acronym},
    name        = {ED},
    long        = {},
    description = {École Doctorale}
}

\newglossaryentry{gpu}
{
    type        = {acronym},
    name        = {GPU},
    long        = {\textbf{G}raphics \textbf{P}rocessing \textbf{U}nit},
    description = {Graphics Processing Unit}
}

\newglossaryentry{hdr}
{
    type        = {acronym},
    name        = {HDR},
    long        = {\textbf{H}abilitation à \textbf{D}iriger des \textbf{R}echerches},
    description = {Habilitation à Diriger des Recherches}
}

\newglossaryentry{ml}
{
    type        = {acronym},
    name        = {ML},
    long        = {\textbf{M}achine \textbf{L}earning},
    description = {Machine Learning}
}

\newglossaryentry{msc}
{
    type        = {acronym},
    name        = {MSc},
    long        = {\textbf{M}aster of \textbf{Sc}ience},
    description = {Master of Science}
}

\newglossaryentry{pdf}
{
    type        = {acronym},
    name        = {PDF},
    long        = {\textbf{P}ortable \textbf{D}ocument \textbf{F}ormat},
    description = {Portable Document Format}
}

\newglossaryentry{phd}
{
    type        = {acronym},
    name        = {PhD},
    long        = {\textbf{Ph}ilosophiæ \textbf{D}octor},
    description = {Philosophiæ Doctor}
}

\newglossaryentry{sota}
{
    type        = {acronym},
    name        = {SOTA},
    long        = {\textbf{S}tate \textbf{O}f \textbf{T}he \textbf{A}rt},
    description = {State Of The Art}
}
